\lecture{10}{Mon 10 Apr 2023 16:33}{Lp space}

\subsection{Norm in Lp Space}

\begin{theorem}[Minkowski Inequality]
	Let $\|f_p\|  = \left( \int _\Omega \left| f \right| ^p d \mu \right)^{\frac{1}{p}}$ .
	where $a \le p \le \infty $, $f \in L(\Omega, \mathcal{A},\mu)$.
	Then
	\[
		\|f + g\|_p \le \|f\|_p + \|g\|_p
	.\] 
\end{theorem}

\begin{definition}[Convexity]
	$\varphi:(a, b) \to \mathbb{R}$ is \textit{convex} if for any $x, y \in (a, b)$, $\lambda \in (0,1)$,
	\[
		\varphi(\lambda x + (1 - \lambda) y) \le  \lambda \varphi(x) + (1 - \lambda) \varphi(y)
	.\] 
\end{definition}

If $\varphi''$ exists, then
\[
	\varphi \text{ convex} \iff \varphi'' > 0
.\] 

\begin{example}
	$\varphi(x) = e^x$ is convex.
	\[
		e^{\lambda x + (1 - \lambda) y} \le  \lambda e^x + (1 - \lambda) e^{y}
	.\] 
	If $e^x = u$, $e^y = v$, then
	\[
		u^\lambda v ^{1 - \lambda} \le  \lambda u + (1 - \lambda) v \quad (u , v > 0)
	.\] 
	Then we have an inequality between weighted arithmetic and geometric means.
\end{example}

\begin{proposition}[Young's Inequality]
	If $p, q>1$, and $\frac{1}{p} + \frac{1}{q} = 1$, then 
	\[
		a b \le \frac{a^p}{p} + \frac{b^q}{q}
	.\] 
	We say $p, q$ are \textit{conjugate exponents}.
\end{proposition}
\begin{proof}
	Set $\lambda = \frac{1}{p}$, then $1 - \lambda = \frac{1}{q}$. Set $a = u^\lambda$ and $b = v^\lambda$. From the inequality above,
	\[
		ab \le  \frac{1}{p} a^p + \frac{1}{q} b^q
	.\] 
	If $a = 0$ or $b=0$, then there is nothing to prove.
\end{proof}

\begin{proposition}[Hölder's Inequality]
	If $f \in L^p$, $g \in L^q$, $p, q > 1$, $\frac{1}{p} + \frac{1}{q } = 1$, then $fg \in L^1$, and
	\[
		\int |fg| \le \left( \int |f|^p \right) ^{\frac{1}{p}} \left( \int |g|^q \right) ^{\frac{1}{q}}
	.\] 
\end{proposition}
Special Case: When $p = q = 2$, this is \logseq{Cauchy-Schwarz Inequality}.

\begin{proof}
	\begin{align*}
		\int |f(x)| |g(x)| d \mu(x)
		&\le \int \frac{|f(x)|^p}{p} + \frac{|g(x)|^q}{q} d \mu(x)
	.\end{align*}
	If $\int |f|^p = \int |g|^q$, then it simplifies to
	\[
		\int |f|^p \left( \frac{1}{p} + \frac{1}{q} \right)  = \int |f|^p  = \left( \int |f|^p \right) ^{\frac{1}{p}} \left( \int |g|^q \right) ^{\frac{1}{q}}
	.\] 
	In general, if $f = 0$ or $g = 0$ a.e., then the inequality is obvious since LHS is zero.

	If both $f$ and $g$ are not zero a.e., then $\int |f|^p$ and $\int |g|^q$ are strictly positive. Let $t$ be such that 
	\[
		\int |tf|^p = \int |g|^p
	.\] 
	Now
	\[
		\int |tf| |g| d \mu \le \left( \int |tf|^p \right) ^{\frac{1}{p}} \left( \int |g|^q \right) ^{\frac{1}{q}}
	.\] 
	Then factor out $t$ and we get the wanted inequality.
\end{proof}

\begin{proof}
	(Proof of Minkowski Inequality) 
	\begin{align*}
		\left( \int |f+g|^p \right) ^{\frac{1}{p}} 
		&\le \left( \int \left( |f| + |g| \right) ^p \right) ^{\frac{1}{p}} 
	.\end{align*}
	Set $h = |f| + |g|$. Then
\begin{align*}
	\int h^p &= \int h^{p-1} |f| + \int h^{p-1}|g| \\
		 &\le \sqrt[q]{\int h^{(p-1) q}} \left( \sqrt[p]{\int |f|^p} + \sqrt[p]{\int |g|^p}   \right)  
.\end{align*}
Where  we have used Hölder's Inequality. Note that
\[
	\frac{1}{q} + \frac{1}{p} = 1 \iff  p = (p-1)q
.\] 
Hence
\begin{align*}
	\int h^p &\le  \left( \int  h^p \right) ^{\frac{1}{q}} \left( \|f\|_p + \|g\|_p \right) \\
\frac{\int h^p}{ \left( \int  h^p \right) ^{\frac{1}{q}}} &\le  \left( \|f\|_p + \|g\|_p \right) \\
 \left( \int  h^p \right) ^{\frac{1}{p}} &\le  \left( \|f\|_p + \|g\|_p \right) \\
 \|h\|_p &\le  \left( \|f\|_p + \|g\|_p \right) 
.\end{align*}

Above claim is true for $\int  h^p \neq 0$. Otherwise,  $h = 0$ a.e., then $f = g = 0$ a.e., and Minkowski inequality trivially follows. 
\end{proof}

\subsection{Lp metric}

Let $\rho(f, g) = \| f - g\|_p$, for $f, g \in L^p$. Is $\rho$ a metric? Symmetry is obvious. Triangle inequality is given by Minkowski Inequality. Nondegeneracy not true in general:
\[
	\rho(f, g) = 0 \implies \int |f - g|^p = 0 \implies f = g \quad a.e.
.\] 
When does $f = g$ a.e. imply $f = g$? Only when the only set of measure zero is $\emptyset $. 

We introduce a equivalence relation of \textbf{a.e. equality}. 

We refer to the quotient space by $L^p(\Omega, \mathcal{A}, \mu)$, the Lp space.

It can be checked that $L^p$ is a vector space.
The Lp norm satisfies
\[
	\rho([f], [g]) = 0 \iff  [f] = [g]
.\] 
where on the LHS the value of $\rho$ does not depend on choice of representatives.

\begin{notation}
	$f \in L^p$ is a ``function'' $\Omega \to [-\infty , \infty ]$, defined almost-everywhere, up to almost-everywhere coincidence.
\end{notation}

\subsection{Generalize to infinity}

\begin{definition}
	We say $f \in L^\infty (\Omega, \mathcal{A}, \mu)$ if $f$ measurable and essentially bounded, i.e. there exists $M \in \mathbb{R}$ such that $|f| \le M$ a.e.

	The least such $M$ is denoted
	\[
		\text{ess} \, \sup  |f| = \|f\|_{\infty }
	.\] 
\end{definition}

We may verify that Minkowski, Holder, and triangle inequality holds when $p = \infty $.

\begin{example}
	For Holder inequality,
	\[
		\int |fg| \le \int \|f\|_\infty = \|f\|_{\infty } \|g\|_1
	.\] 
\end{example}

\subsection{Properties of Lp space}

\begin{theorem}[Riesz-Fisher Theorem]
	$L^p$ is complete. 
\end{theorem}
\begin{proof}
	Let $f_n \in L^p$ be a Cauchy sequence, we need to produce $f$ such that $f_n \to f$ in $L^p$ metric. When $p = \infty $, for all $\varepsilon$ exists $\nu$ such that for all $n, m > \nu$,
	\[
		\text{ess sup} |f_n - f_m| < \varepsilon
	.\] 
	Thus there exists $E_{nm}$, $m(E_{nm}) = 0$, and
	\[
		\text{sup}_{E_{n, m}^c} |f_n - f_m| < \varepsilon
	.\] 
	Let $E = \cup_{n, m > \nu} E_{n, m}$, then $m(E) = 0$. Then on $E^c$, $f_n$ satisfies uniform Cauchy criterion. Hence $f_n$ converges to some $f$ in $E^c$. We can then extend $f$ arbitrarily on $E$. Hence $f_n \to f$ a.e.

	When $p < \infty $, we use
	\begin{lemma}
		Suppose $g_n \in L^p$, $\sum_{n=1}^{\infty} \|g_n\|_p < \infty $, then $\sum_{n=1}^\infty  g_n$ converges a.e. to $g \in L^p$, and 
		\[
			\|g - \sum_{n=1}^{k} g_n\|_p \to  0 \qquad \text{ as }k\to \infty 
		.\] 
		So, $\sum_{n=1}^{\infty} g_n$ also converges in $L^p$.
	\end{lemma}
	\begin{proof}
		We have $\sum_{n=1}^{\infty} \left| g_n \right|  \in L^p$. We consider its $L^p$ norm:
		\begin{align*}
			\int  \left( \sum_{n=1}^{\infty} |g_n| \right) ^p
			&= \lim_k \int \left(\sum_{n=1}^{k} |g_n| \right)^p & \text{MCT}
		.\end{align*}
		Now by Minkowski,
		\begin{align*}
			\|\sum_{n=1}^{k} |g_n|\|_p 
			&\le \sum_{n=1}^{k} \|g_n\|_p \\
			&\le \sum_{n=1}^{\infty} \|g_n\|_p =:S
		.\end{align*}
		Hence
		\[
			\int \left(\sum_{n=1}^{k} |g_n| \right)^p \le S^p
		.\] 
		Therefore
		\[
			\int  \left( \sum_{n=1}^{\infty} |g_n| \right) ^p \le S^p < \infty 
		.\] 
		Hence $ \sum_{n=1}^{\infty} |g_n| $ is finite a.e. Hence  $\sum_{n=1}^{\infty} g_n$ converges a.e., equal to $g$.

		By triangle inequality,
		\[
			|g|^p \le \left( \sum_{n=1}^{\infty} |g_n| \right) ^p \in L^1
		.\] 
		Hence $g^p \in L^1$, hence $g \in L^p$.

		Finally,
		\[
			\|g - \sum_{n=1}^{k} g_n\|_p 
			= \|\sum_{n=k+}^{\infty} g_n\|_p
			\le  \|\sum_{n=k+1}^{\infty} |g_n|\|_p
			\le \sum_{n=k+1}^{\infty} \|g_n\|_p \to 0 \quad \text{as } k \to \infty 
		.\] 
		
	\end{proof}
	\begin{lemma}
		In any metric space, if a Cauchy sequence has a convergent subsequence, then whole sequence converges.
	\end{lemma}

	Suppose $f_n \in L^p$ is Cauchy, we need to find a subsequence that converges. For $k \in \mathbb{N}$ choose $n_k \in \mathbb{N}$ increasing such that
	\[
		\|f_n - f_m\|_p < 2^{-k}
	.\] 
	whenever $n, m \ge  n_k$.

	Define $g_k = f_{n_k} - f_{n_{k-1}}$. Then
	\[
		\sum_{k=1}^{\infty} \|g_k\|_p \le \sum_{k=1}^{\infty} 2^{1-k}
	.\] 
	By lemma 1, 
	\[
		\sum_{k=1}^{m} g_k \to  g \quad \text{a.e. as }m\to \infty  \text{ in }L^p
	.\] 
	Thus
	\[
		f_{n_k} = f_{n_1} + \sum_{j=2}^{k} g_j \to  f_{n_1} + g
	.\] 
	By Lemma 2,
	\[
		f_n \to f_{n_1} + g
	.\] 


	
\end{proof}

Is $L^p$ separable? Not when $p = \infty$. 

\begin{definition}
	$(\Omega, \mathcal{A}, \mu)$ is separable if there exists $\mathcal{A}_0 \subset \mathcal{A}$ countable, with the following property:

	$\forall  E \in \mathcal{A}, \mu(E) < \infty , \forall  \varepsilon>0, \exists A \in \mathcal{A}_0$ such that $\mu(A \Delta E) < \varepsilon$.
\end{definition}

\begin{example}
	$\mathbb{R}$ with Lebesgue measure is separable.

	To see this, let $\mathcal{A}_0$ be generated by all rational bounded intervals. Then $\mathcal{A}_0$ countable.

	Let $\varepsilon>0$, $E$ measurable with $m(E) < \infty $. Choose intervals $I_j$ such that $E \subset \bigcup_{i \in \mathbb{N}} I_j$, and $\sum_{j=1}^{\infty} |I_j| < m(E) + \varepsilon$. WLOG, let $I_j$ be all rational. Choose $N$ such that $\sum_{j=N+1}^{\infty} |I_j| < \varepsilon $.

	Let $A = \bigcup_{j \le N} I_j \in \mathcal{A}_0$.  Then
	\[
		E \setminus A \subset \bigcup_{j \in \mathbb{N}}  I_j \setminus  \bigcup_{j \ge N} I_j
		\subset \bigcup_{j > N} I_j
	.\] 
	Thus $\mu(E \setminus A) < \varepsilon$.

	On the other hand,
	\[
		A \setminus E \subset \bigcup_{j \in \mathbb{N}} I_j \setminus E
	.\] 
	Hence 
	\[
		\mu(A \setminus E)
		< (\mu(E) + \varepsilon) - \mu(E) = \varepsilon
	.\] 
\end{example}

\begin{theorem}
	If $(\Omega, \mathcal{A}, \mu)$ is separable, $1 \le p < \infty $, then $L^p(\Omega, \mathcal{A}, \mu)$ is separable.
\end{theorem}
\begin{proof}
	None.
\end{proof}

\begin{theorem}
	Let $\Omega \subset \mathbb{R}^m$ be measurable, $(\Omega, \mathcal{M}|_\Omega, m|_\Omega)$ is the measure space. If $1 \le p < \infty $, then 
	\[
		C_* (\Omega) = \{f|_\Omega: f \in C(\mathbb{R}^n), \text{vanishes outside some ball}\} 
	.\] 
	is dense in $L^p$.
\end{theorem}
\begin{proof}
	We prove this when $\Omega = \mathbb{R}$.

	We show if $f \in L^p$, $\varepsilon>0$, then there exists some $g \in C(\mathbb{R}), g = 0$ outside some ball, such that $\|f - g\|_p < \varepsilon$.

	First consider the case where $f = \chi_{[a,b]}$. Define $g$ to be
	\[
		g = 
		\begin{cases}
			f \text{ on } (-\infty , a], [a + \frac{\varepsilon^p}{2}, b - \frac{\varepsilon^p}{2}], [b, \infty )\\
			\text{linear in between}
		\end{cases}
	.\] 
	Then
	\[
		\int_{\mathbb{R}}|f - g|^p
		< 
		\int _a^{a + \frac{\varepsilon^p}{2}} 1 + 
		\int _{b - \frac{\varepsilon^p}{2}}^b 1
		= \varepsilon^p
	.\] 
	Hence $\|f - g\|_p < \varepsilon$.

	Next consider $f$ to be linear combination of indicator functions for intervals. Then we can choose $g_j$ such that
	\[
		\|\chi_{[a_j, b_j]} - g_j\|_p < \frac{\varepsilon}{M}
	.\] 
	Then
	\[
		\|f - \sum_{j=1}^{M} c_j g_j\| = \|\sum_{j=1}^{M} c_j(\chi_{[a_j, b_j]} - g_j)\|_p \le \sum_{j=1}^{M} c_j\| \chi_{[a_j, b_j]} - g_j\|_p < \frac{\varepsilon}{M} \sum_{j=1}^{M} c_j
	.\] 

	Now conider general $f$. By HW, for $p = 1$, there exists piecewise constant function $h: \mathbb{R} \to  \mathbb{R}$, such that $\| f - h \|_p < \frac{\varepsilon}{2}$. Then we apply part 2 to each constant component of $h$.
\end{proof}

\begin{proposition}(Translation is continuous)
	Let $1 \le p < \infty $, $f \in L^p(\mathbb{R}^n)$, for  $t \in \mathbb{R}^n$, define
	\[
		f_t(x)  =f(x - t)
	.\] 
	Then $\mathbb{R}^n \ni t \mapsto f_t \in L^p$ is continuous.
\end{proposition}
\begin{proof}
	Apply the density argument. First verify this for $f \in C_0(\mathbb{R}^n)$, continuous and compact supported functions. (HW) Next, for general $f \in L^p$, choose $g_k \in C_0(\mathbb{R}^n)$ such that $g_k \to f$ uniformly in $L^p$. Then $t \mapsto g_{k, t}$ converge to $t \mapsto f_t$ uniformly. 

	Indeed:
	\begin{align*}
		\int _{\mathbb{R}^n} \left| g_{k, t} - f_t \right| ^p 
		&= \int _{\mathbb{R}^n} \left| g_{k}(x - t) - f(x - t) \right| ^p d m(x)\\
		&= \int _{\mathbb{R}^n} \left| g_{k}(x) - f(x) \right| ^p d m(x) \to  0
	.\end{align*} 

	Hence $t \to f_t$ is uniform limit of continuous functions, hence continuous.
\end{proof}

\begin{theorem}
	If  $\Omega\subset R^n$ is bounded, Polynomials aer dense in $L^p$.
\end{theorem}
\begin{proof}
	Use Weierstrass approximation theorem to approximate $f \in C_* (\Omega)$.
\end{proof}

\begin{corollary}
	If $\Omega \subset \mathbb{R}^n	$ is bounded then $L^p(\Omega)$ is separable (with polynomials with coefficients in $\mathbb{Q}$ ).
\end{corollary}


\begin{definition}[Banach Space]
	A Banach space over $\mathbb{R}$ or $\mathbb{C}$ is a vector space $V$ over $\mathbb{R}$ or $\mathbb{C}$, with a norm $\|\cdot\|: V \to [0, \infty )$, such that the metric space induced is complete.
\end{definition}

\begin{example}
	$L^p$ space with $L^p$ norms are Banach spaces.

	If $M$ is a metric space, $C_b(M)$ is continuous bounded functions from $M$ to $\mathbb{R}$, then the $L^\infty $ norm makes $C_b(M)$ a Banach space.
\end{example}

\begin{definition}
	 If $(V, \|\cdot\|_V)$, $(W, \|\cdot\|_W)$ are Banach spaces, then
	 \[
	 	L(V, W) = \{\varphi: V \to  W \text{ linear continuous}\} 
	 .\] 
\end{definition}

\begin{example}
	For $V, W$ euclidean spaces, any linear map is uniquely determined by the image of standard basis $\varphi e_1, \varphi e_2 , \varphi e_m$. Then we can construct matrices.
\end{example}

\begin{definition}
	$V, W$ are \textit{isomorphic} if there exists $\varphi \in L(V, W)$ such that $\varphi$ bijective and $\varphi^{-1} \in L(W, V)$. 

	$V, W$ are isometrically isomorphic if such $\varphi$ also preserves norms.
\end{definition}

\begin{proposition}
	A linear map $\varphi:V \to W$ is continuous iff there exists $C \in \mathbb{R}$ such that for all $v \in V$, $\|\varphi(v)\|_W\le C \|v\|_V$.
\end{proposition}
\begin{proof}
	Necessarity:
	\[
		\|\varphi(v) - \varphi(w)\|_W = \|\varphi(v - w)\|_W \le C \|v - w\|_V
	.\] 

	Sufficiency:
	Since $\varphi$ is continuous at $0 \in V$, there exists $\delta>0$ such that $\varphi(B_\delta(0)) \subset [0,1]$.
	Now for any $v \neq 0$, we can scale it by $\frac{\delta}{\|v\|_V}$ and apply the above.
\end{proof}

\begin{definition}
	For $\varphi \in L(V, W)$, let $\|\varphi\| = \inf \{C \in \mathbb{R}: \forall  v \|\varphi(v)\|\le C \|v\|\} $
\end{definition}
Then, $L(V, W)$ is a vector space.


\begin{theorem}
	$L(V, W)$ is a Banach space.
\end{theorem}

